% ==============================================================================
% REFERENCIAS BIBLIOGRÁFICAS (FORMATO APA)
% ==============================================================================
% Instrucciones de la plantilla:
% - Lista de fuentes de información en formato APA.
% - IMPORTANTE: Enlistar preferentemente libros, revistas, capítulos.
% - Si se citan sitios web, hacerlo conforme al estilo APA.
% ==============================================================================

\renewcommand{\refname}{Referencias Bibliográficas}
\begin{thebibliography}{99}

\bibitem{hearn_baker}
Hearn, D., Baker, M. P., \& Carithers, W. (2011). \textit{Computer Graphics with OpenGL} (4.ª ed.). Pearson Prentice Hall.

\bibitem{foley}
Foley, J. D., Van Dam, A., Feiner, S. K., \& Hughes, J. F. (1996). \textit{Computer Graphics: Principles and Practice in C} (2.ª ed.). Addison-Wesley Professional.

\bibitem{learnopengl}
de Vries, J. (2020). \textit{Learn OpenGL: Learn modern OpenGL graphics programming in a step-by-step fashion}. Recuperado de \url{https://learnopengl.com/}

\bibitem{khronos_opengl}
Khronos Group. (s.f.). \textit{OpenGL 3.3 Reference Pages}. The Khronos Group Inc. Recuperado de \url{https://registry.khronos.org/OpenGL-Refpages/}

\bibitem{teodoro_vite}
Teodoro-Vite, S., et al. (2019). \textit{Fundamentos y Técnicas del Cómputo Gráfico}. Facultad de Ingeniería, Universidad Nacional Autónoma de México.

\end{thebibliography}
